% =====================================================================
%  main_aaai_skeleton.tex  --  AAAI-27 Main Technical Track (MAS)
%  COMPRESSION SCAFFOLD. Target: 7 pages technical + unlimited references.
%
%  HOW TO USE:
%   1. Download the kit from https://aaai.org/authorkit27/ and place
%      aaai27.sty + aaai27.bst in this folder.
%   2. Pour body text from main_standard.tex into the marked slots BELOW,
%      obeying the per-section PAGE BUDGET comments. The budgets sum to 7.0.
%   3. Anything marked [-> APPENDIX] goes in the supplementary technical
%      appendix (unlimited, due 3 days after the paper deadline) and does
%      NOT count toward the 7 pages.
%   4. The hypothesis sentence marked <<HYPOTHESIS>> must appear VERBATIM in
%      three places: abstract, end of intro, conclusion. Do not paraphrase.
%
%  AAAI STYLE HARD RULES (do not violate -> desk reject):
%   - two-column aaai27 style, US Letter, Type1/TrueType fonts
%   - DO NOT change margins, font size, or line spacing
%   - \usepackage of disallowed packages is forbidden; keep the kit preamble
%   - references do NOT count toward 7 pages; technical content does
% =====================================================================

\documentclass[letterpaper]{article}
\usepackage{aaai27}                 % <-- from the downloaded author kit
\usepackage{times}
\usepackage{helvet}
\usepackage{courier}
\usepackage[hyphens]{url}
\usepackage{graphicx}
\graphicspath{{../}{../docs/}}
\usepackage{booktabs}
\usepackage{amsmath,amssymb,amsthm}
\usepackage{newfloat}
\usepackage{caption}
\DeclareCaptionStyle{ruled}{labelfont=normalfont,labelsep=colon,strut=off}
\frenchspacing
\setlength{\pdfpagewidth}{8.5in}
\setlength{\pdfpageheight}{11in}

\theoremstyle{plain}
\newtheorem{theorem}{Theorem}
\newcommand{\wcrit}{w_{\mathrm{crit}}}
\newcommand{\PT}{\mathrm{PT}}

% --- AAAI required metadata commands (fill these in) ---
\pdfinfo{
/Title(Coordination-Collapse Phase Boundary in Policy-Gradient Cooperative MARL)
/Author(Agriya Yadav)
}
\setcounter{secnumdepth}{0}

\title{Coordination-Collapse Phase Boundary in Policy-Gradient Cooperative MARL}
\author{
    %% AAAI requires anonymized submission for the main track --
    %% leave author block per kit instructions for review; de-anonymize for camera-ready.
    Anonymous Submission
}

\begin{document}
\maketitle

% =====================================================================
% ABSTRACT  [BUDGET: included in p.1]
% Keep <= ~200 words. MUST contain the <<HYPOTHESIS>> sentence verbatim.
% Lead the behavioral number with the HUMAN-RANGE lambda~3 (+0.089),
% demote lambda=7 to "high-stakes upper bound".
% =====================================================================
\begin{abstract}
% [POUR: compressed abstract from main_standard lines 56-95]
% Required beats, in order:
%   (1) collapse is a known failure never measured as a quantity;
%   (2) we measure w_crit as a pre-registered, CI'd phase boundary;
%   (3) <<HYPOTHESIS>> sentence verbatim;
%   (4) replication: 3 environments x 2 algorithms;
%   (5) one-line scope honesty: affine law + cascade are SC-specific.
%
% <<HYPOTHESIS>> (verbatim, occurrence 1 of 3):
% "Human-like loss aversion forces cooperative learning agents to require more
%  reward centralization to avoid collapse---a shift we predict in sign a priori
%  and confirm across three environments and two algorithms---and a single
%  downstream loss-averse agent can set that requirement for the whole system."
\end{abstract}

% =====================================================================
% 1. INTRODUCTION  [BUDGET: 1.25 pages]
% Use the 7-point spine from docs/if_this_paper_disappeared.md.
% ORDER CHANGE vs main_standard: lead with the supply-chain justification
% (directed flow -> upstream/downstream definable -> the position question
% is askable). It is currently the 5th reason; make it the 1st.
% End the section with the <<HYPOTHESIS>> sentence (occurrence 2 of 3).
% =====================================================================
\section{Introduction}
\label{sec:intro}
% [POUR: rewritten intro -- see main_aaai_intro.tex when ready]
% CUT FROM main_standard: collapse the four \subsection blocks
% (Coordination collapse / Why PT and MARL / Approach / Why supply chain)
% into 3 tight paragraphs. Move the detailed prior-work contrast to a
% Related Work paragraph or [-> APPENDIX claim ledger].

% =====================================================================
% 2. PRELIMINARIES AND METHOD  [BUDGET: 1.0 page]
% KEEP: CTDE + w dial; CPT value function (Eq cpt); CPT-as-reward-transform
%       (Eq reward) -- this is the anti-"reward-engineering" point, keep it
%       prominent; pre-registered collapse def (Eq collapse); logistic +
%       bootstrap (Eq logistic); three-way decomposition design.
% CUT:  Firth digression -> one sentence + [-> APPENDIX].
% =====================================================================
\section{Preliminaries and Method}
\label{sec:method}
% [POUR: main_standard lines 332-434, compressed]

% =====================================================================
% 3. EXPERIMENTAL SETUP  [BUDGET: 0.5 page]
% KEEP: environment one-paragraph; algorithm one-paragraph; Table experiments.
% CUT:  per-experiment prose -> the table carries it.
% =====================================================================
\section{Experimental Setup}
\label{sec:setup}
% [POUR: main_standard lines 437-488, compressed; keep tab:experiments]

% =====================================================================
% 4. RESULTS  [BUDGET: 1.75 pages -- the core, protect this budget]
%   4.1 Structural transition + metastability  (keep fig:heatmap)
%   4.2 Behavioral decomposition + full lambda sweep
%       -> MERGE tab:baselines INTO tab:decomposition (see MERGED TABLE below)
%       -> keep fig:decomposition; affine-law fig -> can go [-> APPENDIX]
% =====================================================================
\section{Results}
\label{sec:results}
% [POUR: main_standard lines 494-727, compressed]
%
% ---- MERGED TABLE (replaces tab:baselines + tab:decomposition) ----
% One table: each w_crit value, its condition, its ROLE, and the excess vs null.
% This kills kill-shot #2 ("just a fit") by showing roles, not raw midpoints.
%
% \begin{table}[t]\centering\small
% \caption{Behavioral decomposition against the single canonical null.
%   One structural null; every other value is a labeled offset from it.}
% \label{tab:decomp_merged}
% \begin{tabular}{lcccc}
% \toprule
% Condition & $\wcrit$ & 95\% CI & Excess vs null & Role \\
% \midrule
% Null ($\alpha{=}\beta{=}1,\lambda{=}1$), 300k & 0.685 & [0.657,0.710] & 0 & canonical null \\
% \quad reproduced (Exp N)                       & 0.685 & ...           & +0.019 & null, reproduced \\
% \quad at 150k horizon                          & 0.822 & [0.785,0.858] & ---    & null, shorter horizon \\
% Curvature ($\alpha{=}\beta{=}0.88,\lambda{=}1$)& 0.586 & [0.539,0.624] & $-0.100$ & curvature result (eases) \\
% Loss aversion $\lambda{=}3$ (human range)      & 0.774 & [0.751,0.798] & $+0.089$ & behavioral, significant \\
% Loss aversion $\lambda{=}7$ (high-stakes bound)& 0.851 & [0.845,0.858] & $+0.166$ & behavioral, significant \\
% \bottomrule
% \end{tabular}
% \end{table}
%
% NOTE: full lambda=1..7 sweep table (tab:lambda_sweep) -> [-> APPENDIX],
% cite it inline; keep only the 7-point affine-law RMSE 0.014 sentence in body.

% =====================================================================
% 5. MEAN-FIELD MECHANISM  [BUDGET: 1.0 page -- HARD CAP]
% OPEN with the honesty frame, not the theorems:
%   "We claim this mechanism at the sign-and-form level only; magnitude is
%    calibrated, and we report one pre-registered probe that came back
%    inconclusive."
% KEEP IN BODY: Eq drift, Eq snr, Eq wcrit (det. limit), Theorem 3 (sign),
%   the claim ledger table (tab:ledger) -- it is your single best trust object.
% MOVE [-> APPENDIX]: Theorems 1-2 + all proofs; Eq langevin/diffusion detail;
%   Eq sde; the derived affine-law algebra (keep only Eq wcrit_lambda + RMSE).
% =====================================================================
\section{A Mean-Field Mechanism for the Collapse Boundary}
\label{sec:meanfield}
% [POUR: main_standard lines 730-1087, AGGRESSIVELY compressed to 1 page]
% Collapse Thm 1 + Thm 2 into one stated Proposition (one paragraph), proofs out.
% Keep Theorem 3 (sign invariance) IN BODY with a 2-line proof sketch; full
% proof [-> APPENDIX].

% =====================================================================
% 6. GENERALIZATION  [BUDGET: 0.75 page]
% KEEP IN BODY: fig:generalization (the 4-setting summary) + 2 paragraphs:
%   one for "direction generalizes across 3 envs x 2 algos", one for
%   "curvature sign + functional form are substrate-specific (LBF flip)".
% MOVE [-> APPENDIX]: MPE setup detail, LBF setup detail, MAPPO-on-SC detail,
%   fig:mpe_lambda_curves, fig:mappo_curves, the credit-assignment +
%   action-representation ablations (sec:ablations entirely).
% =====================================================================
\section{Generalization across Environments and Algorithms}
\label{sec:generalization}
% [POUR: main_standard lines 1090-1327, compressed to summary + 1 figure]

% =====================================================================
% 7. A MINORITY SETS THE BOUNDARY  [BUDGET: 0.5 page -- KEEP, do not cut]
% This is the most portable, most memorable result. Keep tab:hetero in body.
% Frame as falsifiable structurally-specified prediction (conditions i-iii).
% fig:hetero -> can go [-> APPENDIX] if space is tight; table stays.
% =====================================================================
\section{A Minority Sets the Boundary}
\label{sec:heterogeneous}
% [POUR: main_standard lines 1417-1479, compressed; keep tab:hetero]

% =====================================================================
% 8. DISCUSSION  [BUDGET: 0.5 page]
% KEEP: operational reading (w>0.851 vs w>0.685 cost); lambda~2.25 calibration
%       caveat; the limitations paragraph (substrate scope honesty).
% CUT:  Yamaguchi/Capannoli prose -> one sentence each, fold into Related Work.
% =====================================================================
\section{Discussion}
\label{sec:discussion}
% [POUR: main_standard lines 1482-1584, compressed]

% =====================================================================
% 9. CONCLUSION  [BUDGET: 0.25 page]
% End with the <<HYPOTHESIS>> sentence (occurrence 3 of 3, verbatim).
% =====================================================================
\section{Conclusion}
\label{sec:conclusion}
% [POUR: main_standard lines 1587-1615, compressed]
% <<HYPOTHESIS>> (verbatim, occurrence 3 of 3) goes here.

% =====================================================================
%  PAGE BUDGET LEDGER (must sum <= 7.0 technical pages)
%   Intro            1.25
%   Method           1.00
%   Setup            0.50
%   Results          1.75
%   Mechanism        1.00
%   Generalization   0.75
%   Minority         0.50
%   Discussion       0.50
%   Conclusion       0.25
%   ----------------------
%   TOTAL            7.50  -> 0.5 OVER. Recover by:
%     - moving affine-law figure to appendix (already planned)  ~0.25
%     - tightening Mechanism to 0.85 and Intro to 1.10          ~0.30
%   Re-measure after first pour; trim figures before trimming claims.
%
%  SUPPLEMENTARY TECHNICAL APPENDIX (unlimited, separate file):
%   - Theorems 1-2 + all proofs; full Langevin/SDE derivation
%   - Full lambda-sweep table (tab:lambda_sweep) + affine-law figure
%   - MPE / LBF / MAPPO-on-SC setup + curves
%   - Two ablations (action-representation, credit-assignment)
%   - Threshold robustness; pre-registered SNR probe (report the null!)
%   - Exploration-noise ablation + dissociation
%   - Claim--evidence ledger (tab:claim_ledger); notation table
%   - Reproducibility checklist (REQUIRED by AAAI-27)
% =====================================================================

\bibliographystyle{aaai27}
\bibliography{references_aaai}

\end{document}
